\section{Metodologia Experimental}

\subsection{Plataforma de Simulação e Modelo do Robô}
% TODO: Descrever o ambiente (Ubuntu 24.04, ROS 2 Jazzy Jalisco, Gazebo Sim 8.11.0 Harmonic) e a geometria calibrada do robo diferencial (base_link no centro do eixo motriz, colisoes esfericas pontuais com L_eff = 0,35 m, sensores e limites cinematicos/dinamicos).
% Referenciar a Tabela~\ref{tab:parametros}.

\begin{table}[htbp]
\centering
\caption{Parâmetros do robô diferencial, sensores e controladores avaliados.}
\label{tab:parametros}
\setlength{\tabcolsep}{3.5pt}
\resizebox{\columnwidth}{!}{%
\begin{tabular}{llc}
\toprule
\textbf{Componente} & \textbf{Parâmetro} & \textbf{Valor} \\
\midrule
Robô diferencial & Raio das rodas ($R$) & $0{,}05\text{ m}$ \\
                 & Bitola entre rodas ($L$) & $0{,}35\text{ m}$ \\
                 & Limites cinemáticos ($v_{\max}, \omega_{\max}$) & $0{,}40\text{ m/s}$, $1{,}20\text{ rad/s}$ \\
\midrule
Sensores         & Frequência dos encoders / odometria RK2 & $100\text{ Hz}$ / $50\text{ Hz}$ \\
                 & Ruído nos encoders ($\sigma_\omega$) & $0{,}05\text{ rad/s}$ \\
                 & LiDAR 2D planar ($360$ feixes, $20\text{ Hz}$) & $[0{,}12\text{ m},\, 8{,}0\text{ m}]$ \\
                 & Ruído no LiDAR ($\sigma_\rho$) & $0{,}01\text{ m}$ \\
\midrule
APF (nominal)    & Ganhos ($k_{\text{att}}, k_{\text{rep}}$) & $1{,}0$, $0{,}8$ \\
                 & Raio de influência repulsiva ($\rho_0$) & $1{,}10\text{ m}$ \\
                 & Ganhos cinemáticos ($k_v, k_\omega$) & $0{,}25$, $1{,}8$ \\
\midrule
Bug2             & Distâncias ($d_{\text{hit}}, d_{\text{wall}}, d_{\text{debounce}}$) & $0{,}52\text{ m}$, $0{,}45\text{ m}$, $0{,}45\text{ m}$ \\
                 & Tolerância à linha-$m$ ($\delta_m$) e progresso ($\epsilon$) & $0{,}10\text{ m}$, $0{,}15\text{ m}$ \\
\midrule
Protocolo        & Tolerância à meta ($d_{\text{tol}}$) & $0{,}20\text{ m}$ \\
                 & Sementes aleatórias ($N=5$) & $\{42, 43, 44, 45, 46\}$ \\
\bottomrule
\end{tabular}%
}
\end{table}


\subsection{Protocolo da Etapa I --- Validação da Odometria}
% TODO: Descrever as 3 trajetorias de teste em arena livre (mundo_etapa1_vazio.sdf):
% 1. Linha reta (5 m)
% 2. Arco circular (R_c = 1 m)
% 3. Percurso composto em "S"
% Explicar a frequencia de integracao (50 Hz), uso de use_sim_time e comparacao contra /ground_truth/odom.

\subsection{Protocolo da Etapa II --- Cenários de Navegação e Estudo Paramétrico}
% TODO: Descrever os 3 cenarios (arquivos .sdf), posicoes de partida S=(0,0) e meta T, e o estudo parametrico:
% - Cenario A (cenario_a_esparsos.sdf): obstaculos convexos esparsos + corredor, T=(8,0).
% - Cenario B (cenario_b_armadilha_u.sdf): parede concava em "U", T=(6,0).
% - Subcenario GNRON (cenario_gnron.sdf): parede a 0,45 m atras da meta T=(6,0).
% - Estudo parametrico do APF no Cenario A (Fraco, Nominal, Forte, Ganho Alto) e variacao de rho_0 no GNRON (0,35; 0,80; 1,10; 1,50 m).

\subsection{Métricas de Avaliação e Repetibilidade}
% TODO: Definir formalmente:
% - Etapa I: erros instantaneos e_p(t), e_theta(t) e RMSE_p, RMSE_theta.
% - Etapa II: Taxa de sucesso (%), Colisoes, Comprimento L_traj (m), Tempo total T_tot (s), Distancia minima aos obstaculos d_min (m), Taxa de bloqueio / minimos locais (%), Erro odométrico final e_p(t_f) e Distancia real a meta d_meta (m).
% - Repetibilidade: N = 5 sementes fixas {42, 43, 44, 45, 46} com ruido branco gaussiano nos encoders (sigma_omega = 0,05 rad/s) e no LiDAR (sigma_rho = 0,01 m).
